\section{The frontier}

\begin{theorem}[Parent-barrier gap holds]
  \label{thm:collatz-gap-holds}
  \lean{Infomaterial.Closure.Collatz.parentBarrierGap_holds}
  \notready
  \uses{thm:collatz-gap}
  \emph{(Open frontier --- stated, not proved.)} The gap property itself: nontrivial parent
  barriers never represent $C$. Supported by exhaustive search (1.16M parent barriers, zero
  nonterminal representations) but unproved. The residue refinement is the modular gap
  $D\cdot r-C\ge M$ with $r=C\cdot D^{-1}\bmod M$; stating it needs a modular-inverse primitive
  absent from Lean core (\cref{thm:collatz-residue-gap}).
\end{theorem}
\begin{theorem}[Every nonterminal seed has a crossing]
  \label{thm:collatz-crossing-holds}
  \lean{Infomaterial.Closure.Collatz.everyNonterminalHasCrossing_holds}
  \notready
  \uses{thm:collatz-crossing-firstpassage}
  \emph{(Open frontier --- stated, not proved.)} No nonterminal seed runs an infinite
  coefficient-noncontracting barrier; equivalently every $n>2$ has a finite crossing. This is
  the ``no positive infinite barrier stream'' content.
\end{theorem}
\begin{theorem}[Crossing reduces to "eventually contracts"]
  \label{thm:collatz-crossing-reduces}
  \lean{Infomaterial.Closure.Collatz.everyNonterminalContracts_implies_crossing}
  \leanok
  \uses{thm:collatz-crossing-holds}
  A \emph{first} coefficient crossing exists as soon as \emph{any} contracting prefix does: the
  least contracting length is automatically a first crossing (\texttt{firstCoeffCrossing\_of\_contracting},
  via a from-scratch well-ordering \texttt{exists\_least} since this core lacks \texttt{Nat.find}).
  Hence the sharp form \texttt{EveryNonterminalContracts} --- every $n>2$ has \emph{some}
  contracting parity prefix ($3^{\mathrm{oddCount}}<2^{\mathrm{len}}$) --- \textbf{implies}
  \texttt{EveryNonterminalHasCrossing}, dropping the all-prior-noncontracting bookkeeping. The hole
  is thus reduced to a bare ``every nonterminal seed eventually contracts.''
\end{theorem}
\begin{theorem}[The purl: descent contains the crossing]
  \label{thm:collatz-purl}
  \lean{Infomaterial.Closure.Collatz.everyNonterminalContracts_of_allFallBelow}
  \leanok
  \uses{thm:collatz-descent-iff, thm:collatz-crossing-reduces}
  The surplus $S=A\cdot n+C-B\cdot n$ is the single gauge: descent $\iff S<0$. Its two faces are
  the coefficient crossing ($A<B$) and the offset gap ($C<(B-A)\cdot n$). Because $C\ge0$, $S<0$
  forces $A<B$ (\texttt{contracting\_of\_descent}): a descending step \emph{is} a contracting step.
  Hence universal descent below self contains the crossing hypothesis
  (\texttt{everyNonterminalContracts\_of\_allFallBelow}) --- crossing and gap are the knit and purl
  of one stitch, not two independent holes. (This unifies their \emph{statement}; the gauge's sign
  for every seed remains the open conjecture.)
\end{theorem}
\begin{theorem}[Accelerated Collatz]
  \label{thm:collatz-accelerated}
  \lean{Infomaterial.Closure.Collatz.accelerated_collatz}
  \notready
  \uses{thm:collatz-conditional, thm:collatz-gap-holds, thm:collatz-crossing-holds}
  \emph{(Open frontier --- derived, rests on the two holes.)} Every positive seed reaches the
  terminal cycle. The derivation from \cref{thm:collatz-conditional} is complete; the node is
  red because it depends on exactly \cref{thm:collatz-gap-holds} and
  \cref{thm:collatz-crossing-holds} and nothing else.
\end{theorem}
